\chapter{Regression and Applied Statistics}

A small café records nine trading days. Two numbers are written every evening: the number of
online orders in tens, and the kitchen's overtime hours. The café owner needs one number for
one question: if orders rise by ten, how many extra overtime hours must she prepare for
tomorrow?

\begin{center}
\small
\begin{tabular}{l*{9}{r}}
\hline
day                & 1 & 2 & 3 & 4 & 5 & 6 & 7 & 8 & 9 \\
orders (tens)      & 1 & 2 & 3 & 4 & 5 & 6 & 7 & 8 & 9 \\
overtime (hours)   & 10 & 12 & 13 & 14 & 16 & 18 & 19 & 20 & 22 \\
\hline
\end{tabular}
\end{center}

A quick answer is available without any tool: from the first day to the last, orders rise by
eight tens and overtime rises by twelve hours, so roughly $1.5$ hours per ten orders. Least
squares, which we will build in this chapter, gives $1.466667$. The two numbers are close, and
precisely because they are close we are tempted to treat the matter as settled. Three
important questions are left unanswered by both. How far does that number move if those nine days
are replaced by nine other days? If tomorrow's orders are eighty, how many hours must be
prepared? And if one day turns out to be misrecorded, how much damage does that do?

Chapter 10 already gave the picture: least squares is the orthogonal projection of the
observation vector onto the column space of the design matrix, and the normal equations say
that the residual is perpendicular to that space. What is missing there is uncertainty.
Chapter 20 already gave confidence intervals and hypothesis tests, but for a single parameter
from observations that are independent and share one distribution. Here the parameter sought
is a slope, a ratio of two sums that are both random, and the observations are no longer on
equal footing: a point at the end of the range carries more influence than a point in the
middle. This chapter joins the two and stops at an honest place: how far a conclusion may be
carried.

\begin{untukanak}
A grandmother sells coconut ice in front of her house. She never records anything neatly, but
her memory is good enough: on a hot day her stock sells out before the afternoon. On a cloudy
day half the ice melts in the cooler. One day her grandchild arrives with an exercise book and
records nine days in a row: roughly how hot it was, and how many cups were sold.

From those nine records the grandchild draws one straight line: the hotter it gets, the more is
sold. The grandmother nods, but then asks, ``if tomorrow is a little hotter, how many cups
should I prepare?'' The grandchild falls silent. The line gives the direction, not the number.
Besides, on one day it rained hard although it was hot, and the sales were dreadful; that point
sits far from the line.

In this chapter we will learn why the best-fitting line is still not a sure forecast, why one
strange day can shift the line further than five ordinary days, and why anyone who sells
numbers must state how wide their uncertainty is.
\end{untukanak}

\section{Nine days and one line}
\tingkat{basic}

Nine pairs of numbers, two unknowns, nine equations. We met a system like this in Chapter 8:
more equations than unknowns, so an exact solution exists only if all nine points happen to be
collinear. No line passes exactly through every point above, and that is not a measurement
failure but the usual state of affairs. All we can look for is the closest line, and
``closest'' has to be defined first.

In everyday language: we look for the one line whose total error is as small as possible, with
the error counted as the perpendicular distance from the point to the line.

\kamus{regression}{regresi}{drawing the one line closest to all the points, rather than a line
that passes through all of them}

\begin{definisi}[Simple linear regression model]
\index{regression}
For observations $i=1,\dots,n$,
\[
y_i = a + b\,x_i + \varepsilon_i ,
\]
with $x_i$ the observed explanatory variable, $y_i$ the explained variable, $a$ the intercept,
$b$ the slope, and $\varepsilon_i$ the unobserved error. The model is called \emph{linear}
because it is linear in the parameters $a$ and $b$, not because the line is straight in $x$.
\end{definisi}

One small distinction governs the whole chapter: $\varepsilon_i$ and the residual are not the
same object. The error $\varepsilon_i$ is the difference between the true value and the
population line that we never see; the residual $e_i=y_i-\hat y_i$ is the difference between
the observed value and the line we fitted from the data itself. A residual can be computed and
drawn, an error can only be assumed. Whenever a book states an assumption about errors and then
checks it through a picture of residuals, there is one jump that must not be forgotten: what is
inspected is a stand-in, not the thing itself.

\section{Least squares as a projection}
\tingkat{intermediate}

Stack $y\in\mathbb{R}^n$ as the observation vector and let $X$ be the $n\times2$ matrix whose
columns are one and $x$. Finding the best line means choosing $\hat y$ in the column space of
$X$ that is closest to $y$: the problem of Chapter 10, with the addition that we need the
explicit form of the coefficients in order to compute their uncertainty.

\begin{definisi}[Least-squares estimator]
\index{least squares}
$\hat\beta=(\hat a,\hat b)$ is the least-squares estimator if it minimizes
\[
\operatorname{RSS}(\beta)=\|y-X\beta\|^2=\sum_{i=1}^{n}\bigl(y_i-a-bx_i\bigr)^2 .
\]
\end{definisi}

\kamus{least squares}{kuadrat terkecil}{adding up all the errors after squaring them, then
choosing the line whose total is smallest}

\begin{teorema}[The normal equations]
\index{normal equations}
$\hat\beta$ minimizes RSS if and only if
\[
X^{\top}(y-X\hat\beta)=0 ,
\qquad\text{that is}\qquad X^{\top}X\hat\beta=X^{\top}y .
\]
If the columns of $X$ are linearly independent, $\hat\beta=(X^{\top}X)^{-1}X^{\top}y$ is
unique.
\end{teorema}

\begin{proof}[Proof]
Write $f(\beta)=\|y-X\beta\|^2=y^{\top}y-2\beta^{\top}X^{\top}y+\beta^{\top}X^{\top}X\beta$.
Its gradient is $-2X^{\top}y+2X^{\top}X\beta$, so the stationary condition is exactly the normal
equations. Since $X^{\top}X$ is positive definite when its columns are linearly independent,
$f$ is strictly convex and that stationary point is the global minimum.

The converse direction is more informative. Suppose $X^{\top}(y-X\hat\beta)=0$. For any
$\beta$,
\[
\|y-X\beta\|^2=\|y-X\hat\beta\|^2+\|X(\hat\beta-\beta)\|^2
+2(\hat\beta-\beta)^{\top}X^{\top}(y-X\hat\beta),
\]
and the cross term vanishes. Hence $\operatorname{RSS}(\beta)\ge\operatorname{RSS}(\hat\beta)$
for every $\beta$, with equality only when $X\beta=X\hat\beta$.
\end{proof}

That proof contains the whole geometric content: the residual is perpendicular to every column
of $X$. Since the first column is the vector of ones and the second is $x$, we get
$\sum_i e_i=0$ and $\sum_i x_ie_i=0$, and those two identities are what support the
decomposition of sums of squares in the next section.

\begin{proposisi}[Hat matrix]
\index{hat matrix}
With $H=X(X^{\top}X)^{-1}X^{\top}$ we have $\hat y=Hy$, $H^{\top}=H$, $H^2=H$,
$\operatorname{tr}H=2$, and $0\le h_{ii}\le1$ with $\sum_i h_{ii}=2$. The value $h_{ii}$ is
called the \emph{leverage} of point $i$.
\end{proposisi}

\begin{proof}[Proof]
Symmetry and idempotence follow directly from the definition:
\[
H^2=X(X^{\top}X)^{-1}X^{\top}X(X^{\top}X)^{-1}X^{\top}=H .
\]
Hence $H$ is an orthogonal projection matrix, so
$h_{ii}=e_i^{\top}He_i=\|He_i\|^2\ge0$ and $h_{ii}\le\|e_i\|^2=1$. Finally
$\operatorname{tr}H=\operatorname{tr}\bigl((X^{\top}X)^{-1}X^{\top}X\bigr)
=\operatorname{tr}I_2=2$.
\end{proof}

Now the numbers, because the whole chapter rests on one computation that can be rechecked. For
the café data,
\[
X^{\top}X=\begin{pmatrix}9&45\\45&285\end{pmatrix},
\qquad X^{\top}y=\begin{pmatrix}144\\808\end{pmatrix},
\qquad \det(X^{\top}X)=9\cdot285-45^2=540 .
\]
Therefore
\[
\hat\beta=\frac1{540}\begin{pmatrix}285&-45\\-45&9\end{pmatrix}
\begin{pmatrix}144\\808\end{pmatrix}
=\frac1{540}\begin{pmatrix}4680\\792\end{pmatrix}
=\begin{pmatrix}8.666667\\1.466667\end{pmatrix}.
\]
The number $808$ is worth checking by hand: $\sum x_iy_i=10+24+39+56+80+108+133+160+198$. So is
$4680$: $285\cdot144-45\cdot808=41040-36360$.

One more theorem is needed before we may speak about the quality of an estimator. That theorem
explains why least squares is used in almost every field, and at the same time why it is not a
sacred choice.

\begin{teorema}[Gauss--Markov]
\index{gauss-markov}
Assume (A1) $y=X\beta+\varepsilon$; (A2) $\mathbb{E}[\varepsilon\mid X]=0$;
(A3) $\operatorname{Var}(\varepsilon\mid X)=\sigma^2I$; (A4) the columns of $X$ are linearly
independent. Then $\hat\beta$ is unbiased with $\operatorname{Var}(\hat\beta)=\sigma^2
(X^{\top}X)^{-1}$, and among all linear unbiased estimators, $\hat\beta$ has the smallest
variance. Such an estimator is called \emph{BLUE} (\emph{best linear unbiased estimator}).
\end{teorema}

\begin{proof}[Proof]
Unbiasedness: $\mathbb{E}[\hat\beta\mid X]=(X^{\top}X)^{-1}X^{\top}X\beta=\beta$, using (A2).
Variance: $\hat\beta-\beta=(X^{\top}X)^{-1}X^{\top}\varepsilon$, so
$\operatorname{Var}(\hat\beta)=(X^{\top}X)^{-1}X^{\top}(\sigma^2I)X(X^{\top}X)^{-1}
=\sigma^2(X^{\top}X)^{-1}$.

The third part is not proved in full here. The sketch: every linear unbiased estimator has the
form $\tilde\beta=Cy$ with $CX=I$, so $\tilde\beta-\beta=C\varepsilon$ and
$\operatorname{Var}(\tilde\beta)=\sigma^2CC^{\top}$. Writing
$C=(X^{\top}X)^{-1}X^{\top}+D$ gives $DX=0$ and
\[
CC^{\top}-(X^{\top}X)^{-1}=DD^{\top},
\]
which is positive semi-definite, so every component of $\tilde\beta$ has variance no smaller
than the matching component of $\hat\beta$. The full proof is in Seber and Lee,
\emph{Linear Regression Analysis}, Chapter 3, and in Greene, \emph{Econometric Analysis},
Chapter 3.
\end{proof}

\section{Coefficients and their units}
\tingkat{basic}

The slope $1.466667$ means nothing until its unit is named. The unit is hours per ten orders:
every additional ten orders raises the forecast overtime by $1.466667$ hours, that is $88$
minutes, or $8.8$ minutes per order. The café owner can use that sentence; the sentence ``the
coefficient is $1.47$'' she cannot.

The intercept $8.666667$ hours is the forecast for $x=0$, a day with no online orders at all.
No such day occurs in the data, and outside the range $1$ to $9$ the line is not guaranteed to
hold. Calling that number ``fixed overtime cost'' is \index{extrapolation}extrapolation wrapped
in a name: we are using the line outside the place where it was fitted. What is safe: the
intercept is the number that makes the line pass through $(\bar x,\bar y)$, because
$\hat a=\bar y-\hat b\bar x$; it is useful for computation, not for interpretation.

A change of units does not change the conclusion. If $x$ is written in orders rather than tens,
the slope is $0.146667$ hours per order and its standard error is $0.004714$, both divided by
ten; $t$ and $R^2$ do not move. Multiplying one variable by a constant $c$ multiplies the
coefficient and its standard error by $c$ or $1/c$, while their ratio stays put, so a report
that quotes $t$ without units cannot be read.

\begin{tangan}
Prepare graph paper, nine coins, and a piece of string as long as the side of the paper. Place
the coins on the grid according to the nine pairs of numbers at the start of the chapter:
horizontally orders, vertically overtime, one square to one unit. Tie the ends of the string to
the two edges of the paper so that the string can be slid as a straight line. Slide and rotate
the string until the sum of the perpendicular distances from all the coins to the string feels
smallest; do not measure, just compare by eye, then lock it with tape.

What should be visible: two people will put the string in slightly different places, and the
difference is largest when you shift the ends of the string. Now lift the leftmost coin by one
square and notice how far the string has to turn. Put it back, then lift the middle coin by one
square: almost nothing changes. That is leverage, and you have just seen it with your hands
before meeting the formula.
\end{tangan}

\section{R squared and its correlation}
\tingkat{intermediate}

Once the line exists, people ask ``how good is it''. The most common answer is one number that
compares the residual with the original spread of $y$. In everyday language: what share of the
wobble in $y$ the line manages to explain, and what share is left over.

\begin{definisi}[$R^2$]
\index{coefficient of determination}
With $\mathrm{SST}=\sum_i(y_i-\bar y)^2$, $\mathrm{SSE}=\sum_i e_i^2$, and
$\mathrm{SSR}=\sum_i(\hat y_i-\bar y)^2$, the coefficient of determination is
\[
R^2=1-\frac{\mathrm{SSE}}{\mathrm{SST}}=\frac{\mathrm{SSR}}{\mathrm{SST}} .
\]
\end{definisi}

\kamus{coefficient of determination}{koefisien determinasi}{the share of the wobble that the
line manages to explain, expressed as a fraction of the whole}

\begin{teorema}[Decomposition of sums of squares]
If $X$ contains a column of ones, then $\mathrm{SST}=\mathrm{SSR}+\mathrm{SSE}$ and
$0\le R^2\le1$.
\end{teorema}

\begin{proof}[Proof]
Write $y-\bar y\mathbf{1}=(\hat y-\bar y\mathbf{1})+e$. Since $\sum e_i=0$, the vector $e$ is
perpendicular to $\bar y\mathbf{1}$, and by the normal equations $e$ is also perpendicular to
$\hat y$. Hence the two terms on the right are perpendicular, and Pythagoras gives
$\|y-\bar y\mathbf{1}\|^2=\|\hat y-\bar y\mathbf{1}\|^2+\|e\|^2$.
\end{proof}

For the café data, $\mathrm{SST}=130$, $\mathrm{SSE}=0.933333$, so
$R^2=1-0.933333/130=0.992821$. In a simple regression, $R^2$ equals the squared correlation:
$r=88/\sqrt{60\cdot130}=0.996404$ and $r^2=0.992821$. That equality is special to one
explanatory variable; as soon as there are two, $R^2$ is no longer the square of any
correlation.

The number $0.992821$ easily makes people stop thinking, so let us tame it. What is explained is
$99.28$ percent of the \emph{wobble} in overtime, and that wobble is measured against its own
mean. The residuals still have variance $s^2=0.933333/7=0.133333$, that is a standard deviation
$s=0.365148$ hours or $21.9$ minutes. A café owner who prepares overtime on the basis of this
line can still miss by more than twenty minutes. $R^2$ is relative; $s$ is absolute, and what
people feel is $s$.

\begin{untukanak}
Imagine two students. The first runs a hundred-meter race and finishes $0.2$ seconds behind the
winner. The second runs a marathon and finishes twenty minutes behind the winner. Both lose, but
not in the same sense. Two tenths of a second is very close; twenty minutes is very far. If the
teacher writes only ``second place'' for both, we lose the whole story.

Numbers of that kind occur in this chapter. There is a number that measures the \emph{share}
guessed correctly, and there is a number that measures the \emph{size} of the error left over.
The first can look almost perfect while the second stays large. Someone who reads only the first
number will prepare too little, then wonder why the forecast misses.
\end{untukanak}

\begin{proposisi}[Adding a variable never lowers $R^2$]
If a column block $X_2$ is added to $X_1$ so that its column space contains the column space of
$X_1$, then $\mathrm{SSE}$ does not increase and $R^2$ does not decrease.
\end{proposisi}

\begin{proof}[Proof]
The set of vectors $X_1\beta_1$ is contained in the set $X_2\beta_2$, so the minimum distance to
$y$ taken over the larger set cannot be larger. Hence $R^2$ cannot fall, even when the added
variable has nothing to do with $y$.
\end{proof}

The practical consequence matters: comparing models by raw $R^2$ always favors the fatter model.
The usual correction is $R^2_{\text{adj}}=1-(1-R^2)(n-1)/(n-p)$ with $p$ the number of
parameters. For the multiple model of the next section, $R^2=0.997951$ falls to $0.997496$ when
$p=3$ and $n=12$; at small $n$ with many variables the correction can reverse the ranking.

\section{Multiple regression and collinearity}
\tingkat{advanced}

So far, one explanatory variable. In reality almost every problem demands more than one, and
that is where its characteristic disease appears: two explanatory variables carrying almost the
same information.

Twelve branches of a food stall are recorded: daily sales $y$ in millions of rupiah, the number
of employees $x_1$, and the seating area $x_2$ in square meters.

\begin{center}
\small
\begin{tabular}{l*{12}{r}}
\hline
branch   & 1 & 2 & 3 & 4 & 5 & 6 & 7 & 8 & 9 & 10 & 11 & 12 \\
$x_1$    & 6 & 7 & 7 & 8 & 8 & 9 & 9 & 10 & 10 & 11 & 12 & 12 \\
$x_2$    & 18 & 22 & 20 & 26 & 24 & 28 & 26 & 32 & 30 & 34 & 36 & 38 \\
$y$      & 11.2 & 12.4 & 12.0 & 14.1 & 13.6 & 15.2 & 14.8 & 17.0 & 16.4 & 18.1 & 19.4 & 19.9 \\
\hline
\end{tabular}
\end{center}

Least squares gives
\[
\hat y=2.442991+0.686916\,x_1+0.239252\,x_2,
\]
with standard errors $0.207583$, $0.135366$, and $0.042101$; the $t$ values are
$11.7687$, $5.0745$, and $5.6828$ on $9$ degrees of freedom; $s=0.143607$ and
$R^2=0.997951$. All of it looks healthy. Then we fit one explanatory variable at a time: $y$ on
$x_1$ gives $1.446019$, $y$ on $x_2$ gives $0.450075$, both nearly twice their multiple
coefficients.

Not because of a miscalculation. Employees and floor area almost determine each other: their
correlation is $r_{12}=0.986801$. Once $x_2$ is in, $x_1$ carries almost no new information, so
its coefficient shrinks to its additional effect only, namely $0.687$ million per employee at
fixed area. The single coefficient $1.446$ mixes the employee effect with the area effect that
usually accompanies it, and the gap between the two is a measure of how far the two variables
overlap.

\begin{definisi}[Variance inflation factor]
\index{variance inflation}
For explanatory variable $j$, let $R_j^2$ be the coefficient of determination of the regression
of $x_j$ on all the other explanatory variables. Its variance inflation factor is
\[
\mathrm{VIF}_j=\frac{1}{1-R_j^2}.
\]
\end{definisi}

\kamus{variance inflation factor}{varians inflasi}{how much the standard error of a coefficient
is blown up by other explanatory variables that carry similar information}

\begin{proposisi}[VIF with two explanatory variables]
\index{collinearity}
If there are only two explanatory variables, $\operatorname{Var}(\hat b_1)=\sigma^2/\bigl(S_{11}
(1-r_{12}^2)\bigr)$, with $S_{11}=\sum_i(x_{i1}-\bar x_1)^2$. Hence the standard error of
$\hat b_1$ is inflated by exactly $\sqrt{\mathrm{VIF}_1}$ compared with the case of mutually
orthogonal explanatory variables.
\end{proposisi}

\begin{proof}[Proof]
For two explanatory variables, inverting the $2\times2$ matrix of the centered block
$X^{\top}X$ gives first element $S_{22}/(S_{11}S_{22}-S_{12}^2)$, and
$S_{12}=r_{12}\sqrt{S_{11}S_{22}}$. Dividing numerator and denominator by $S_{22}$ yields
$\sigma^2/\bigl(S_{11}(1-r_{12}^2)\bigr)$.
\end{proof}

The numbers can be checked: $S_{11}=42.9167$, $r_{12}^2=0.973775$, so $\mathrm{VIF}=38.132232$
and $\sqrt{\mathrm{VIF}}=6.1751$. The standard error of $\hat b_1$ is indeed $0.135366$, while
$\sigma/\sqrt{S_{11}}$ is only $0.021921$; the ratio is $6.1751$, matching the root of the VIF
to four digits.

Collinearity does not damage the forecast; it damages the interpretation. The value
$R^2=0.997951$ shows that the line is still close to the points, and a forecast at $x_1=9$,
$x_2=28$ is still useful. What cannot be trusted is the sentence ``one extra employee raises
sales by $0.687$ million'': its $95$ percent interval $(0.380697;\ 0.993135)$ is wide, and if
the data shift a little that coefficient can change a lot without the line changing much.

\section{Intervals, the $t$ test, and prediction intervals}
\tingkat{advanced}

The next question is the same as in Chapter 20, only the object changes: not a mean, but a
slope. We need one quantity whose distribution is known when the model is correct, then we
invert it into an interval.

In everyday language: we measure how far the observed slope sits from zero, in units of the
slope's own typical wobble.

\kamus{prediction interval}{selang prediksi}{an interval for one day that has not happened yet,
always wider than the interval for the mean because it has to bear the surprise of that day as
well}

\begin{teorema}[The $t$ test and an interval for one coefficient]
Assume (A1) to (A4) and normal errors. With $p$ the number of parameters,
$s^2=\mathrm{SSE}/(n-p)$, and $\operatorname{se}(\hat b_j)=s\sqrt{[(X^{\top}X)^{-1}]_{jj}}$,
we have
\[
\frac{\hat b_j-\beta_j}{\operatorname{se}(\hat b_j)}\sim t_{n-p},
\]
so the $1-\alpha$ interval for $\beta_j$ is
$\hat b_j\pm t_{n-p,1-\alpha/2}\operatorname{se}(\hat b_j)$, and the test of
$H_0\colon\beta_j=0$ rejects at level $\alpha$ when $|t|>t_{n-p,1-\alpha/2}$.
\end{teorema}

\begin{proof}[Proof]
$\hat b_j$ is a linear combination of $y$ with fixed coefficients, so it is normal with
expectation $\beta_j$ and variance $\sigma^2[(X^{\top}X)^{-1}]_{jj}$ by Gauss--Markov. Moreover
$\mathrm{SSE}/\sigma^2\sim\chi^2_{n-p}$ and is independent of $\hat\beta$, because the residual
is perpendicular to the column space while $\hat\beta$ lies entirely inside it. The ratio of a
standard normal variable to the root of a chi-square variable divided by its degrees of freedom
follows a $t$ distribution. The full proof is in Wasserman, \emph{All of Statistics},
Chapter 13.
\end{proof}

For the café data, $\hat b=1.466667$, $s=0.365148$, $S_{xx}=60$, so
$\operatorname{se}(\hat b)=0.047140$. With $t_{7,0.975}=2.364624$,
\[
\hat b\pm t\,\operatorname{se}(\hat b)=1.466667\pm0.111470
=(1.355197;\ 1.578136),
\]
and $t=31.1127$ with a p-value below $10^{-6}$. Rejecting $\beta=0$ here is no achievement: with
residuals as tight as $22$ minutes and an $x$ range this wide, almost any data would reject it.
What deserves reporting is the interval, because it says that the extra overtime per ten orders
lies between $1.36$ and $1.58$ hours.

Now distinguish two questions that are often confused. What is the average overtime on a day
with eighty orders? And what is the overtime on one particular day with eighty orders? The first
answer narrows like $1/\sqrt n$; the second never narrows below the standard deviation of the
residuals, because a single day brings its own surprise.

\begin{proposisi}[Two kinds of forecast variance]
For fixed $x_0$ and a new observation $y_0=x_0^{\top}\beta+\varepsilon_0$ independent of the old
data,
\[
\operatorname{Var}(\hat y_0)=\sigma^2x_0^{\top}(X^{\top}X)^{-1}x_0,
\qquad
\operatorname{Var}(y_0-\hat y_0)=\sigma^2\bigl(1+x_0^{\top}(X^{\top}X)^{-1}x_0\bigr).
\]
\end{proposisi}

\begin{proof}[Proof]
The first already appeared in Gauss--Markov. For the second, write
$y_0-\hat y_0=\varepsilon_0-x_0^{\top}(\hat\beta-\beta)$; the two terms are independent because
$\varepsilon_0$ is independent of the old data, so the variance is the sum of the two
variances.
\end{proof}

For $x_0=8$ in the café data, $x_0^{\top}(X^{\top}X)^{-1}x_0=1/9+9/60=0.261111$. The half-width
of the mean interval is $2.364624\sqrt{0.133333\cdot0.261111}=0.441209$, so
$(19.9588;\ 20.8412)$ with width $0.882$. The half-width of the prediction interval is
$2.364624\sqrt{0.133333\cdot1.261111}=0.969635$, so $(19.4304;\ 21.3696)$ with width $1.939$.
The ratio of widths is $\sqrt{1.261111/0.261111}=2.198$: the prediction interval is more than
twice as wide, and no amount of data can remove that excess. Figure~\ref{fig:dua-selang} shows
both bands over the whole range of $x$.

\begin{figure}[htbp]
\centering
\begin{tikzpicture}
\begin{axis}[
  width=11.4cm, height=6.0cm,
  xmin=0.6, xmax=9.4, ymin=8.2, ymax=25.4,
  axis lines=left,
  xtick={1,2,3,4,5,6,7,8,9},
  tick label style={font=\scriptsize},
  xlabel={online orders (tens)}, ylabel={overtime hours},
  xlabel style={font=\small}, ylabel style={font=\small},
  legend style={font=\scriptsize,draw=layergrey,at={(0.02,0.98)},
    anchor=north west,row sep=0.5pt},
]
\addplot[accent,very thick,domain=0.7:9.3,samples=60]
  {8.666667+1.466667*x+0.863439*sqrt(1.111111+(x-5)^2/60)};
\addplot[accent,very thick,domain=0.7:9.3,samples=60]
  {8.666667+1.466667*x-0.863439*sqrt(1.111111+(x-5)^2/60)};
\addplot[accent,very thick,domain=0.7:9.3,samples=60,forget plot]
  {8.666667+1.466667*x};
\addplot[layerindigo,densely dashed,thick,domain=0.7:9.3,samples=60]
  {8.666667+1.466667*x+0.863439*sqrt(0.111111+(x-5)^2/60)};
\addplot[layerindigo,densely dashed,thick,domain=0.7:9.3,samples=60]
  {8.666667+1.466667*x-0.863439*sqrt(0.111111+(x-5)^2/60)};
\addplot[only marks,mark=*,mark size=1.8pt,layerindigo] coordinates {
  (1,10) (2,12) (3,13) (4,14) (5,16) (6,18) (7,19) (8,20) (9,22)};
\legend{prediction bands, , fitted line, mean bands, }
\end{axis}
\end{tikzpicture}
\caption{The $95$ percent confidence interval for mean overtime (dashed line) and the $95$
percent prediction interval for one new day (thick line), both around the fitted line
$\hat y=8.666667+1.466667x$. At $x=8$ the mean interval
$(19.9588;\ 20.8412)$ has width $0.882$, while the prediction interval
$(19.4304;\ 21.3696)$ has width $1.939$. Both bands widen at the ends of the range, where the
line is less supported by data; the ratio of their widths, $2.198$, does not depend on $n$.}
\label{fig:dua-selang}
\end{figure}

\section{Residual diagnostics}
\tingkat{intermediate}

Every interval above is conditional on the model being right. Hence the first step after fitting
a line is not to compute p-values but to inspect the residuals. In everyday language: we look at
whether what is left over from the line still has a pattern; a residual with a pattern means the
line has the wrong shape, not that chance is at work.

\index{residual}
Four patterns deserve recognition: a curve against $x$, sign of a non-linear relation; a funnel
that widens with $x$, a sign that (A3) fails; one faraway point on its own, a suspect
misrecording; and a pattern by time order, a sign that observations are not independent.

For the café data, the residuals in order are
\[
-0.1333,\ 0.4,\ -0.0667,\ -0.5333,\ 0,\ 0.5333,\ 0.0667,\ -0.4,\ 0.1333 .
\]
There is no evident curvature, and the sum is zero as it should be. What is interesting is their
size: the largest residual is $0.5333$ hours, about $32$ minutes, on the fourth and sixth days,
both near the center of the data. Those points are in fact \emph{not} very influential, because
their leverage is only $0.1278$.

The leverage of point $i$ in a simple regression has a closed form
\[
h_{ii}=\frac1n+\frac{(x_i-\bar x)^2}{S_{xx}},
\]
so on this data $h_{11}=h_{99}=1/9+16/60=0.3778$, while $h_{55}=0.1111$. An end point holds
$3.4$ times the weight of a middle point, and that is not a designer's choice but a consequence
of geometry: a line rotated about the center of the data moves furthest at the ends. Actual
influence combines the size of the residual with leverage, and the usual measure is Cook's
distance
\[
D_i=\frac{e_i^2}{p\,s^2}\cdot\frac{h_{ii}}{(1-h_{ii})^2},
\]
which here is largest, $0.2870$, on the second and eighth days, both with residual $0.4$ and
leverage $0.2611$. Values below one are generally considered not worrying, but that threshold
is a convention, not a law.

One warning that is rarely written: with $n=9$ diagnostics have almost no power. A pattern in
nine points easily arises from chance, and a pattern that is really there is easily invisible.
Diagnostics offer suspicion, not verdict.

\begin{bacakeras}
Two voices.

``The $R$ squared is zero point nine nine. So the model is almost perfect.''

``Perfect for what? The residuals are still twenty-two minutes.''

``But ninety-nine percent of the wobble is already explained.''

``Explained by what? By orders. Not by a kitchen that suddenly slowed down.''

``So a high $R$ squared is not a guarantee.''

``A high $R$ squared only means the line is tight to the points that exist. It knows nothing
about a day that has not come yet.''

Say it once more, slowly: tight to the old points is not a promise about a new day. Repeat it
until it feels natural in the mouth: tight to the old points, not a promise.
\end{bacakeras}

\begin{resep}
\begin{enumerate}
  \item Plot $y$ against $x$ before computing anything, and ask whether a straight line makes
        sense.
  \item Fit the line by least squares, then write the equation complete with the unit of every
        coefficient.
  \item Inspect the residuals against $x$ and against time order; look for curvature, a funnel,
        and isolated points.
  \item Report $s$ together with $R^2$, then an interval for the coefficient that carries the
        main conclusion.
  \item If the aim is to forecast one event, use a prediction interval; if the aim is a mean,
        use a mean interval.
\end{enumerate}
A warning about order: step 1 must not be swapped with step 2. Plotting after fitting the line
makes the eye look for justification, not for faults.
\end{resep}

\section{When the answer is yes or no}
\tingkat{advanced}

So far $y$ has been a number that can be measured. Many important questions are binary: a
customer returns or not, a transaction is fraudulent or not, a patient recovers or not. It is
still tempting to fit an ordinary straight line to $0$ and $1$, then read it as a probability.
Two things go wrong at once, and both can be computed.

The following data: $26$ customers, the distance of their home from the stall $x$ in kilometers,
and $y=1$ if they returned within thirty days. Of the $26$, $15$ returned. A linear regression
on the binary $y$, called the \emph{linear probability model}, gives
$\hat y=1.346154-0.096154x$. At $x=2$ the forecast is $1.153846$: a probability greater than
one, and that is not a rounding error but a defect of form. At $x=20$, still outside the data
range but only seven kilometers outside, the forecast is $-0.576923$.

The second failure is subtler. For binary $y$, the variance of the residual is really $p(1-p)$
and therefore depends on $x$: at $x=2$ its value is $0.007225$, at $x=14$ it is $0.028890$. The
linear probability model assumes one variance for all points, namely $s^2=0.124199$ on this
data, so its standard errors are wrong exactly where it matters most.
Ordinary least squares, which works well when the variance is constant, here relies on an
assumption that plainly does not hold.

\index{logistic regression}
\begin{definisi}[Logistic model]
For $p_i=P(y_i=1\mid x_i)$,
\[
\log\frac{p_i}{1-p_i}=a+b\,x_i,
\qquad\text{that is}\qquad
p_i=\frac{1}{1+e^{-(a+bx_i)}} .
\]
The parameters are estimated by maximum likelihood, not by least squares, so there is no closed
form and the computation is iterative.
\end{definisi}

\kamus{odds ratio}{rasio peluang}{how many times the odds multiply for one additional unit, not
by how many points the probability rises}

For that data, maximum likelihood gives $\hat a=6.315432$ and $\hat b=-0.699945$ with standard
errors $2.383413$ and $0.258594$. The odds ratio is $e^{-0.699945}=0.496613$: each extra
kilometer cuts the odds of returning to about a half, in the sense of a ratio, not a difference
in points. The estimated probabilities are $0.992722$ at $x=2$ and $0.029777$ at $x=14$, both
legitimate as probabilities. The $95$ percent interval for $b$ is
$(-1.233657;\ -0.166233)$; it is entirely negative, so the direction is clear even if the
magnitude is not.

One costly misreading: a logistic coefficient is not the marginal effect on probability. The
value $b=-0.699945$ acts on the log-odds scale, and its effect on probability is largest around
$p=0.5$ and smallest at the ends. A good report states probabilities at several meaningful
values of $x$, not just the coefficient.

\section{Correlation is not causation}
\tingkat{basic}

Now the part most often violated in reports that look tidy. Two groups of employees attend
training of different length, and the number of errors per hundred transactions is recorded.

\begin{center}
\small
\begin{tabular}{l*{6}{r}}
\hline
morning shift & training hours & 1 & 2 & 3 & 4 & 5 \\
              & errors         & 4.0 & 3.6 & 3.2 & 2.9 & 2.5 \\
\hline
night shift   & training hours & 5 & 6 & 7 & 8 & 9 \\
              & errors         & 8.0 & 7.5 & 7.0 & 6.4 & 6.0 \\
\hline
\end{tabular}
\end{center}

Within the morning shift, the slope is $-0.36$ errors per training hour. Within the night shift,
$-0.474286$. Both are negative: longer training goes with fewer errors. But when the two groups
are pooled, the slope is $+0.359036$ with $R^2=0.226121$. The sign flips, and not one of the
numbers is wrong.

The cause is not magic but the addition of two sources:
\[
b_{\text{pooled}}=\frac{\sum_g S_{xy,g}+\sum_g n_g(\bar x_g-\bar x)(\bar y_g-\bar y)}
{\sum_g S_{xx,g}+\sum_g n_g(\bar x_g-\bar x)^2}.
\]
The morning group contributes $S_{xy}=-6.3$ and $S_{xx}=17.5$; the night group $S_{xy}=-8.3$ and
$S_{xx}=17.5$; the between-group term is $6\cdot(-2)(-1.85)+6\cdot(2)(1.85)=44.4$; the
denominator is $83$. Hence
$b_{\text{pooled}}=(-6.3-8.3+44.4)/83=29.8/83=0.359036$. The between-group term $44.4$ defeats
the sum of within-group terms $-14.6$, and that is the whole secret.

The interpretation: training hours are not assigned at random. Night-shift employees, who make
more errors because of working hours and the kinds of transactions, also receive more training,
so what is measured as the effect of training is mostly the difference between shifts.
Figure~\ref{fig:simpson} places everything on one plane.

\begin{figure}[htbp]
\centering
\begin{tikzpicture}
\begin{axis}[
  width=11.4cm, height=6.2cm,
  xmin=0.5, xmax=10.5, ymin=1.6, ymax=8.9,
  axis lines=left,
  xtick={1,2,3,4,5,6,7,8,9,10},
  tick label style={font=\scriptsize},
  xlabel={training hours}, ylabel={errors per 100 transactions},
  xlabel style={font=\small}, ylabel style={font=\small},
  legend style={font=\scriptsize,draw=layergrey,at={(0.98,0.98)},
    anchor=north east,row sep=0.5pt},
]
\addplot[only marks,mark=*,mark size=2pt,layerteal] coordinates {
  (1,4.0) (2,3.6) (3,3.2) (4,2.9) (5,2.5)};
\addplot[only marks,mark=square*,mark size=1.9pt,layerplum] coordinates {
  (6,7.5) (7,7.0) (8,6.4) (9,6.0) (10,5.7)};
\addplot[layerteal,very thick,domain=0.8:5.4,samples=2] {4.326667-0.36*x};
\addplot[layerplum,very thick,domain=5.6:10.2,samples=2]
  {10.32381-0.474286*x};
\addplot[accent,thick,densely dashed,domain=0.8:10.2,samples=2]
  {2.941968+0.359036*x};
\legend{morning, night, pooled}
\end{axis}
\end{tikzpicture}
\caption{Simpson's paradox in the training data. Within each shift the line falls: $-0.36$ errors
per hour for the morning and $-0.4743$ for the night. The pooled line, computed from all twelve
points regardless of shift, instead rises with slope $+0.3590$. The change of sign comes from
the difference between groups, not from one odd point; the between-group term $44.4$ defeats the
within-group term $-14.6$.}
\label{fig:simpson}
\end{figure}

\begin{untukanak}
Two classes take part in a reading contest. Class A reads thin books, many pages, and gets good
marks. Class B reads thick books, fewer pages, and gets poor marks. If all the pupils are pooled,
it looks as if the fewer pages read, the higher the mark. Yet inside each class, the more pages,
the higher the mark.

What is wrong is not the numbers but the pooling. Class B really was given thicker books, and
class B really is weaker at reading. The two facts blend into one deceptive line.

In this chapter we will learn why pooling two groups can reverse the direction of a conclusion,
why the question ``does training reduce errors'' cannot be answered just by plotting all the
points, and why anyone changing policy must know which groups they are really comparing.
\end{untukanak}

\begin{tantangan}
Build a table with two groups, five points each, such that both within-group slopes are positive
but the pooled slope is negative, and the difference between the group means of $y$ is as small
as still possible. How small is that difference, and why can it not be smaller?

Hint 1: the sign of the pooled slope is decided by the between-group term. Make the first group
have small $x$ and small $y$, the second large $x$ and large $y$, then reverse the direction
inside each group.

Hint 2: with five points and within-group slope $+1$ at $S_{xx}=10$, each group contributes
$S_{xy}=+10$ and the between-group term must be less than $-20$. Since
$n_g(\bar x_g-\bar x)(\bar y_g-\bar y)$ is negative only when one of the mean differences has
the opposite sign to the other, the mean difference must pass the point where the two terms
cancel.

If you get stuck, go back to the section Correlation is not causation and recompute
$29.8/83$ step by step.
\end{tantangan}

\section{Honest reporting}
\tingkat{intermediate}

A regression analysis can be done correctly and reported so as to mislead without a single number
being falsified. Four habits close most of that gap.

First, state how many observations and how many parameters, because the reader needs both to
compute the degrees of freedom and to judge whether $s$ can be trusted. On the café data, $n=9$
and $p=2$, so $s$ is built from seven degrees of freedom; hiding that number makes $s$ look
sturdier than it is.

Second, separate the data used to choose the model from the data used to judge it. Choosing the
model form, choosing the variables, and computing the error on the same set always produces
numbers that are too good. At $n=9$ that separation is impossible, and precisely for that reason
the conclusion must be called weak, not reinforced with $R^2=0.992821$.

Third, state how many models were tried. If ten possibilities are tested and the best is
reported, the chance that at least one of them looks meaningful is $1-0.95^{10}=0.401263$ even
though none of the effects is real. That number is not the probability of being wrong, but the
probability of picking something that is a coincidence.

Fourth, report magnitudes with units and intervals, not just stars. The sentence
``$\hat b=0.687$ ($95$ percent: $0.381$ to $0.993$ million per employee)'' can be checked; the
sentence ``significantly affects'' cannot. If the conclusion is about causation, say why:
randomization, or adjustment for the named confounders.

\section{A worked example from end to end}
\tingkat{intermediate}

We repeat the whole chain on the café data, with the intermediate numbers written out, so that
every step can be checked with a calculator.

The means are $\bar x=5$ and $\bar y=16$. The squared deviations are
$S_{xx}=\sum(x_i-5)^2=60$ and $S_{yy}=130$; their cross product is
$S_{xy}=\sum(x_i-5)(y_i-16)=808-9\cdot5\cdot16=88$. Hence $\hat b=88/60=1.466667$ and
$\hat a=16-1.466667\cdot5=8.666667$.

The forecasts at the nine points are $10.1333$; $11.6$; $13.0667$; $14.5333$; $16$; $17.4667$;
$18.9333$; $20.4$; $21.8667$. The residuals are $-0.1333$; $0.4$; $-0.0667$; $-0.5333$; $0$;
$0.5333$; $0.0667$; $-0.4$; $0.1333$, and their sum of squares is $\mathrm{SSE}=0.933333$. With
$n-p=7$, $s^2=0.133333$ and $s=0.365148$ hours.

The standard error of the slope is $s/\sqrt{S_{xx}}=\sqrt{0.133333/60}=\sqrt{0.00222222}=0.047140$.
The statistic is $t=1.466667/0.047140=31.1127$, and $t_{7,0.975}=2.364624$, so the half-width of
the interval is $0.111470$ and the interval is $(1.355197;\ 1.578136)$.

\kertas{Repeat the same computation for one change of data. The ninth day is recorded with $22$
hours of overtime; suppose the correct record is $26$ hours. Recompute $\hat b$, $s$,
$\operatorname{se}(\hat b)$, and its $95$ percent interval, then compare with the table above.
The correct answers: $\hat b=1.733333$, $s=1.306880$, $\operatorname{se}(\hat b)=0.168718$,
interval $(1.3344;\ 2.1323)$. The width is $0.798$ against $0.223$, that is $3.6$ times.
Explain why one point at the end can do that, when it is only one of nine.}

\section{Computational lab}
\tingkat{advanced}

Every number in this chapter can be recomputed in one short file. The following program uses
$X^{\top}X$ directly, so it also checks the algebra of the second section.

\begin{lab}
\begin{verbatim}
import numpy as np
from math import sqrt
from scipy import stats

x = np.arange(1, 10, dtype=float)
y = np.array([10, 12, 13, 14, 16, 18, 19, 20, 22], dtype=float)
n = len(y)
X = np.column_stack([np.ones(n), x])

beta = np.linalg.solve(X.T @ X, X.T @ y)
fit = X @ beta
e = y - fit
s2 = (e ** 2).sum() / (n - 2)
sst = ((y - y.mean()) ** 2).sum()

print("X'X =", (X.T @ X).tolist(), " X'y =", (X.T @ y).tolist())
print("beta = %.6f %.6f" % tuple(beta))
sse = (e ** 2).sum()
print("SSE = %.6f  R2 = %.6f  s = %.6f"
      % (sse, 1 - sse / sst, sqrt(s2)))
print("residuals =", np.round(e, 4).tolist())

cov = s2 * np.linalg.inv(X.T @ X)
se = np.sqrt(np.diag(cov))
tq = stats.t.ppf(0.975, n - 2)
print("se = %.6f %.6f  t = %.4f %.4f"
      % (se[0], se[1], beta[0] / se[0], beta[1] / se[1]))
lo, hi = beta[1] - tq * se[1], beta[1] + tq * se[1]
print("CI95 slope = (%.6f, %.6f)" % (lo, hi))

for x0 in (5.0, 8.0):
    c = np.array([1.0, x0])
    y0 = c @ beta
    v = s2 * (c @ np.linalg.inv(X.T @ X) @ c)
    print("x0=%.0f  yhat=%.4f  mean +/-%.4f  prediction +/-%.4f"
          % (x0, y0, tq * sqrt(v), tq * sqrt(v + s2)))

h = np.einsum("ij,jk,ik->i", X, np.linalg.inv(X.T @ X), X)
cook = e ** 2 / (2 * s2) * h / (1 - h) ** 2
print("h_ii =", np.round(h, 4).tolist())
print("D_Cook =", np.round(cook, 4).tolist())

x1 = np.array([6, 7, 7, 8, 8, 9, 9, 10, 10, 11, 12, 12], dtype=float)
x2 = np.array([18, 22, 20, 26, 24, 28, 26, 32, 30, 34, 36, 38], dtype=float)
yb = np.array([11.2, 12.4, 12.0, 14.1, 13.6, 15.2, 14.8, 17.0, 16.4, 18.1,
               19.4, 19.9])
Xb = np.column_stack([np.ones(12), x1, x2])
bb = np.linalg.solve(Xb.T @ Xb, Xb.T @ yb)
eb = yb - Xb @ bb
s2b = (eb ** 2).sum() / 9
seb = np.sqrt(np.diag(s2b * np.linalg.inv(Xb.T @ Xb)))
r12 = np.corrcoef(x1, x2)[0, 1]
print("multiple beta = %.4f %.4f %.4f  s = %.4f"
      % (bb[0], bb[1], bb[2], sqrt(s2b)))
print("multiple se   = %.4f %.4f %.4f" % tuple(seb))
print("r12 = %.4f  VIF = %.2f  sqrt VIF = %.2f"
      % (r12, 1 / (1 - r12 ** 2), 1 / sqrt(1 - r12 ** 2)))


def ols(xx, yy):
    A = np.column_stack([np.ones(len(xx)), xx])
    return np.linalg.solve(A.T @ A, A.T @ yy)[1]


xp = np.array([1, 2, 3, 4, 5, 6], dtype=float)
yp = np.array([4.0, 3.6, 3.2, 2.9, 2.5, 2.2])
xm = np.array([5, 6, 7, 8, 9, 10], dtype=float)
ym = np.array([8.0, 7.5, 7.0, 6.4, 6.0, 5.7])
print("Simpson: morning %.4f  night %.4f  pooled %.4f"
      % (ols(xp, yp), ols(xm, ym), ols(np.r_[xp, xm], np.r_[yp, ym])))

xd = np.array([2, 2, 3, 3, 4, 4, 5, 5, 6, 6, 7, 7, 8, 8, 9, 9, 10, 10,
               11, 11, 12, 12, 13, 13, 14, 14], dtype=float)
yd = np.array([1, 1, 1, 1, 1, 1, 1, 1, 1, 1, 1, 0, 1, 0, 1, 0, 1, 0, 0,
               1, 0, 0, 0, 0, 0, 0], dtype=float)
A = np.column_stack([np.ones(len(yd)), xd])
lpm = np.linalg.solve(A.T @ A, A.T @ yd)
print("LPM: %.4f %+.4f  prediction at x=2: %.4f  at x=20: %.4f"
      % (lpm[0], lpm[1], lpm @ [1, 2], lpm @ [1, 20]))


def nll(th):
    z = A @ th
    return float(np.sum(np.log1p(np.exp(-z)) + (1 - yd) * z))


g = np.array([0.0, 0.0])
for _ in range(40):
    p = 1 / (1 + np.exp(-(A @ g)))
    grad = A.T @ (p - yd)
    Hm = A.T @ (A * (p * (1 - p))[:, None])
    g = g - np.linalg.solve(Hm, grad)
se_g = np.sqrt(np.diag(np.linalg.inv(Hm)))
p = 1 / (1 + np.exp(-(A @ g)))
print("logit: %.4f %+.4f  se %.4f %.4f" % (g[0], g[1], se_g[0], se_g[1]))
print("deviance = %.4f  p(x=2) = %.4f  p(x=14) = %.4f"
      % (2 * nll(g), p[0], p[-1]))
print("odds ratio per km = %.4f" % np.exp(g[1]))
\end{verbatim}
\end{lab}

The output, exactly as it comes:

\begin{verbatim}
X'X = [[9.0, 45.0], [45.0, 285.0]]  X'y = [144.0, 808.0]
beta = 8.666667 1.466667
SSE = 0.933333  R2 = 0.992821  s = 0.365148
residuals = [-0.1333, 0.4, -0.0667, -0.5333, 0.0, 0.5333, 0.0667, -0.4, 0.1333]
se = 0.265274 0.047140  t = 32.6706 31.1127
CI95 slope = (1.355197, 1.578136)
x0=5  yhat=16.0000  mean +/-0.2878  prediction +/-0.9101
x0=8  yhat=20.4000  mean +/-0.4412  prediction +/-0.9696
h_ii = [0.3778, 0.2611, 0.1778, 0.1278, 0.1111, 0.1278, 0.1778, 0.2611, 0.3778]
D_Cook = [0.0651, 0.287, 0.0044, 0.1792, 0.0, 0.1792, 0.0044, 0.287, 0.0651]
multiple beta = 2.4430 0.6869 0.2393  s = 0.1436
multiple se   = 0.2076 0.1354 0.0421
r12 = 0.9868  VIF = 38.13  sqrt VIF = 6.18
Simpson: morning -0.3600  night -0.4743  pooled 0.3590
LPM: 1.3462 -0.0962  prediction at x=2: 1.1538  at x=20: -0.5769
logit: 6.3154 -0.6999  se 2.3834 0.2586
deviance = 17.8843  p(x=2) = 0.9927  p(x=14) = 0.0298
odds ratio per km = 0.4966
\end{verbatim}

Two things must be read from that output. First, $\mathrm{X}^{\top}\mathrm{X}$ and
$\mathrm{X}^{\top}y$ are printed exactly as computed by hand, and $\hat\beta$ comes out of
solving the system; agreement of this kind is not proof that the theory is true, but it removes
arithmetic slips, which are the most frequent cause. Second, the last three parts show why
interpretation is more fragile than computation: $0.6869$ against $1.4460$ in the single
regression, VIF $38.13$, and a pooled slope whose sign reverses. Not one of the three is a
programming error.

\begin{untukdosen}
The teaching order I suggest: start from the table of nine days, ask the class to guess the slope
in two ways, by the difference of the ends and by ``the closest line''. The two answers differ
only slightly, and that is exactly the entry point: the real question is not the number but its
wobble. Only after the class is restless, move to the normal equations and show that
$\sum x_ie_i=0$ is a geometric statement, not an algebraic trick.

The two misconceptions that appear most often. First, students read $R^2$ as ``the probability
that the forecast is right'', then are surprised when the prediction interval is two hours wide
at $R^2=0.99$; the symptom shows when they cannot explain why $s$ does not shrink along with
$R^2$. Second, they interpret multiple coefficients like single ones, even though on the
twelve-branch data the number moves from $1.446$ to $0.687$ only because another explanatory
variable enters.

A blackboard activity that works: have two groups compute the slope on the training data, one
group using all twelve points, one separating the shifts. Write both results side by side, then
ask each group to defend theirs. The sign that the class is ready to move on: they can name one
situation in which a multiple coefficient changes sign, and can explain why the forecast stays
good.
\end{untukdosen}

\section{Summary}
\tingkat{basic}

\begin{ringkasan}
\begin{itemize}
  \item Regression fits the line that minimizes the sum of squared residuals. By the normal
        equations the residual is perpendicular to every column of $X$; on the café data
        $\hat b=88/60=1.466667$ and $\hat a=8.666667$.
  \item A coefficient is meaningful only together with its unit: $1.466667$ hours per ten orders
        equals $8.8$ minutes per order. Changing the scale of $x$ changes the coefficient and its
        standard error, not $t$ or $R^2$.
  \item $R^2$ is relative, $s$ is absolute. On the café data $R^2=0.992821$ but $s=0.365148$
        hours, that is $21.9$ minutes.
  \item Under (A1)--(A4) Gauss--Markov guarantees that least squares is BLUE with
        $\operatorname{Var}(\hat\beta)=\sigma^2(X^{\top}X)^{-1}$; that guarantee lapses when the
        error variance is not constant or the explanatory variables are collinear.
  \item Collinearity does not damage the forecast but it damages the interpretation: on the
        twelve branches, $r_{12}=0.986801$ gives $\mathrm{VIF}=38.13$ and standard errors $6.18$
        times larger, while $R^2$ stays at $0.997951$.
  \item A prediction interval is always wider than a mean interval; at $x_0=8$ the ratio is
        $2.198$ and does not shrink with $n$. The sign of a slope can reverse when two groups are
        pooled: $-0.36$ and $-0.474286$ within groups, $+0.359036$ once pooled.
\end{itemize}
\end{ringkasan}

\section{Exercises}
\tingkat{intermediate}

\begin{latihan}
\begin{enumerate}
  \item Evaluate this claim: ``Because $R^2=0.99$, this line can forecast a single day with an
        error of at most twenty minutes.'' True, false, or not enough evidence? Name the part
        that is right, the part that is wrong, and the additional number that settles the
        assessment.
  \item Prove that with a column of ones in $X$ we have $\sum_i e_i=0$ and
        $\sum_i x_ie_i=0$, then explain why both fail when the line is forced through the
        origin.
  \item Prove $\operatorname{Var}(\hat b)=\sigma^2/S_{xx}$ from
        $\hat b=\sum_i(x_i-\bar x)y_i/S_{xx}$ without matrix formulas.
  \item Compute $\hat a$, $\hat b$, $s$, $\operatorname{se}(\hat b)$, and the $95$ percent
        interval for $(1,2)$, $(2,2.5)$, $(3,4)$, $(4,3.5)$, $(5,5)$, then check with Python. Why
        is the interval much wider than the interval for the café data?
  \item Compute the $95\%$ prediction interval for $x_0=2$ and $x_0=9$ on the café data.
        Why do the widths differ when $n$ is the same, and why does the interval never
        shrink to zero when $n$ is increased?
  \item Show that adding explanatory variables cannot raise SSE, then explain why
        $R^2_{\text{adj}}$ can fall. Compute $R^2_{\text{adj}}$ for the café data if $p$ rises
        from $2$ to $3$ at the same $R^2$.
  \item For $S_{11}=50$, $S_{22}=200$, $r_{12}=0.9$, and $\sigma^2=2$, compute the VIF and the
        standard error of $\hat b_1$ with and without $x_2$ in the model.
  \item Build a table with two groups with within-group slope $-0.2$ and pooled slope $+0.4$,
        then check the result with the decomposition formula.
  \item Rewrite this sentence so that it is honest: ``the best model of the twelve we tried gives
        $p=0.03$, so the effect is real''.
\end{enumerate}

Hints for the three hardest items. Item 3: write
$\hat b-\beta=\sum_i(x_i-\bar x)\varepsilon_i/S_{xx}$, take the variance, use independence and
$\operatorname{Var}(\varepsilon_i)=\sigma^2$. Item 6: compare the column spaces of the two
models, then remember the divisor $n-p$ that shrinks. Item 7: $\mathrm{VIF}=1/(1-0.81)=5.263158$,
so the standard error is $\sqrt{5.263158}$ times larger.
\end{latihan}

\begin{untukpeneliti}
What is genuinely proved here: the normal equations as a necessary and sufficient condition for
the minimum, together with uniqueness when $X$ has full column rank; the properties of the hat
matrix and $\operatorname{tr}H=p$; unbiasedness and $\operatorname{Var}(\hat\beta)=\sigma^2
(X^{\top}X)^{-1}$; the Gauss--Markov inequality within the class of linear unbiased estimators;
the decomposition $\mathrm{SST}=\mathrm{SSR}+\mathrm{SSE}$ when a column of ones is present; the
non-monotonicity of $R^2$ under added columns; the $t$ distribution for the coefficient ratio
under normal errors; and the variance of the forecast error.

What is merely popular: that $R^2$ measures the quality of a model, that multiple coefficients
always represent effects, and that the thresholds $\mathrm{VIF}>10$ or $D_i>1$ are derived
results. Primary sources: Seber and Lee, \emph{Linear Regression Analysis}, Chapter 3; Greene,
\emph{Econometric Analysis}, Chapter 3; Wasserman, \emph{All of Statistics}, Chapter 13;
Simpson (1951) for the reversal of direction in contingency tables; Hosmer, Lemeshow, and
Sturdivant, \emph{Applied Logistic Regression}, Chapters 1--3.

Limits of validity: every interval is conditional on the model being right, the errors being
independent, and the variance being constant. On grouped data, time series, or repeated
observations, independence fails and the reported standard errors are too small. Open question:
how far post-selection intervals depart from their nominal level, and how influence diagnostics
are constructed when $n$ is far smaller than $p$.
\end{untukpeneliti}

\peta{Chapter 21 in Wasserman, \emph{All of Statistics}: Chapter 13 (Linear and Logistic
Regression: least squares, the $t$ distribution for coefficients, and logistic regression) and
Chapter 6 (Models, Statistical Inference, and Learning: the interpretation of intervals and
standard errors). The geometry of least squares as an orthogonal projection is set out in
Chapter 10 of this book, where the normal equations first appear as a consequence of
orthogonality, and is used again in Chapter 24 for regularization. Confidence intervals,
p-values, and the interpretation of levels are taken from Chapter 20 and not repeated here. The
Gauss--Markov theorem and its conditions are referred to Seber and Lee Chapter 3 and Greene
Chapter 3, because Wasserman does not prove it. Simpson's paradox is referred to Simpson's
(1951) paper and used again in Chapter 25 when discussing data separation and evaluation
leakage, and in Chapter 26 when discussing causal claims that must be proved, not inferred from
correlation.}
